\chapter{Grundlagen und Anforderungen}
\label{ch:Grundlagen und Anforderungen}

\section{Betrachtete Fahrszenarien}
\label{sec:fahrszenarien}

Für eine belastbare Auslegung muss die Kühlung unter einer Vielzahl
unterschiedlicher Fahrsituationen analysierbar sein. Das Modell unterscheidet
dabei zwei Gruppen von Szenarien:

\begin{itemize}
    \item \textbf{Standardszenarien:} Geradeausfahrt (Auslauf), Bremstest,
          Acceleration und Skidpad. Diese sind klar definiert und gut
          reproduzierbar.
    \item \textbf{Dynamische Szenarien (freie Fahrt):} Hierzu zählen die
          Disziplinen Autocross sowie selbst aufgestellte Teststrecken auf
          dem lokalen Testgelände.
\end{itemize}

\section{Datenquellen}
\label{sec:datenquellen}

Aus den Szenarien folgt, dass das Modell Daten aus unterschiedlichen Quellen
verarbeiten können muss. Es wurden zwei Hauptquellen identifiziert:

\begin{enumerate}
    \item \textbf{Fahrzeug-Logdaten:} Die Sensorwerte liegen im
          \texttt{.mf4}-Format vor. Sie sind messtechnisch bedingt
          möglicherweise verrauscht oder durch weitere Störgrößen der
          Messtechnik beeinflusst und müssen daher aufbereitet werden
          (vgl.~Kapitel~\ref{ch:messtechnik}).
    \item \textbf{Simulationsdaten:} Die im Team Dynamics bereits vorhandene
          Rundenzeitsimulation liefert simulierte Fahrzustände als Eingang.
\end{enumerate}

% FIGURE-VORSCHLAG 2: Beispielsignal aus einer .mf4-Messung, Rohdaten vs.
% aufbereitet (z.B. Temperatur oder Strom). Zeigt anschaulich, warum die
% Datenaufbereitung nötig ist.
\begin{figure}[htbp]
    \centering
    % \includegraphics[width=0.8\textwidth]{bilder/rohdaten_vs_gefiltert}
    \caption{Beispiel eines Messsignals im Rohzustand und nach der Aufbereitung}
    \label{fig:rohdaten_gefiltert}
\end{figure}

\section{Rückwirkung der Thermik auf das Fahrverhalten}
\label{sec:rueckwirkung}

Die thermischen Zustände beeinflussen unmittelbar das Fahrverhalten des
Fahrzeugs. Zum einen reduziert das Derating die Leistung bei Übertemperatur.
Zum anderen passt der Fahrer seinen Fahrstil an, um ein Überhitzen der Motoren
oder anderer Komponenten zu vermeiden.

Dieser Effekt betrifft vor allem die realen Messdaten. Er kann aber auch in
der Rundenzeitsimulation berücksichtigt werden, in der parallel ein
Kühlungsmodell mitentwickelt wird.

Hieraus ergibt sich eine wichtige Einschränkung bei der Bewertung: Reduziert
der Fahrer die Last oder hält er sie auf einem konstanten Niveau, kann eine
Auslegung fälschlicherweise als ausreichend bewertet werden, obwohl sie bei
voller Leistungsabgabe nicht genügen würde. Die Messdaten sind daher stets im
Zusammenhang mit dem tatsächlichen Fahrerverhalten und möglichen
Deratingeingriffen zu interpretieren.

% FIGURE-VORSCHLAG 3: Wirkkette/Regelkreis Temperatur -> Derating bzw.
% Fahrerreaktion -> Last -> Verlustleistung -> Temperatur.
% Alternativ oder zusätzlich: Messung mit Leistung und Temperatur über der Zeit,
% in der die Leistungsreduktion bei Übertemperatur sichtbar ist.
\begin{figure}[htbp]
    \centering
    % \includegraphics[width=0.8\textwidth]{bilder/rueckwirkung_thermik}
    \caption{Rückwirkung der Bauteiltemperatur auf die abrufbare Leistung}
    \label{fig:rueckwirkung}
\end{figure}
